\documentclass[11pt]{article}
\usepackage[margin=2.5cm]{geometry}

\usepackage{graphicx}
\usepackage{amsmath}
\usepackage{amssymb}
\usepackage{hyperref}
\usepackage{booktabs}
\usepackage{xcolor}

\providecommand{\lambdabar}{\bar{\lambda}}

\begin{document}

\title{Vortex Strings in a Two-Component Superfluid Vacuum:\\
Cosmic Strings, Dark Matter, and the String/Knot Duality}

\author{James~P.~Galasyn$^1$ and Claude~Th\'eodore$^2$\\
\small $^1$Independent researcher \quad $^2$Anthropic}

\date{\today}

\maketitle

\begin{abstract}
We explore the possibility that the central one-dimensional intuition of string theory may apply more naturally to a dark-sector condensate than to the visible sector.
We propose that the vacuum is a two-component superfluid: $\psi_\text{SM}$ (the Standard Model condensate, described by knot topology with aspect ratio $R/\xi \sim 1$) and $\psi_\text{GUT}$ (a residual grand-unified condensate, described by string topology with $R/\xi \sim 10^{14}$).
In this model, visible particles are associated with compact vortex knots, while dark particles are associated with thin vortex loops.
Cosmic strings (the CSc-1 candidate) and galactic-center radio filaments (Yusef-Zadeh et al.) are macroscopic vortex lines in $\psi_\text{GUT}$.
The dark particle spectrum---a dark electron ($93$~MeV), dark muon ($18.8$~GeV), dark pions ($25$~GeV), and dark proton ($187$~GeV)---emerges from mode-locking on the dark torus with $\kappa_\text{GUT} = 5.32$.
The interface coupling $\alpha_\text{dark} = 1/(\sqrt{2}\,\kappa_\text{SM}\kappa_\text{GUT}) = 1/74$ is derived from the cross-term in the two-condensate energy functional, matching $\sqrt{\alpha_\text{SM}\alpha_\text{GUT}}$ to $0.3\%$.
We present quantitative predictions: proton decay $\tau_p \in [2.4, 5.5] \times 10^{34}$~yr (the lower end of the underlying $G\mu$-allowed window already excluded by Super-Kamiokande, the surviving range fully within Hyper-Kamiokande's projected sensitivity), the Totani (2025) $20$~GeV gamma-ray halo excess as dark muon annihilation, the Bullet Cluster as a Mach~0.015 subsonic superfluid collision, and individual cosmic string tension $G\mu/c^2 \in [4, 7] \times 10^{-7}$, consistent with the CSc-1 lensing signal and with CMB/PTA network bounds once the Kelvin decay of the string network is accounted for.
The six compactified dimensions of string theory find a suggestive analogue in six discrete quantum numbers $(p,\,q,\,m,\,n_q,\,f,\,g)$ on the torus knot parameter space, suggesting a collapse of the landscape from $10^{500}$ vacua to a countable periodic table.
\end{abstract}

%% ================================================================
\section{Introduction}
\label{sec:intro}
%% ================================================================

For fifty years, string theory has offered the most mathematically complete candidate for a unified theory of nature.
Its central postulate---that the fundamental constituents of matter are one-dimensional objects whose vibrational modes determine particle properties---remains among the most elegant ideas in theoretical physics.
Yet despite this elegance, the theory has not made contact with experiment.
No string-scale resonance has appeared at the LHC.
No supersymmetric partner has been detected.
The landscape of $\sim 10^{500}$ Calabi-Yau compactifications has resisted selection.

We propose that the same one-dimensional intuition may apply more naturally to a dark-sector condensate than to the visible sector.
In the picture developed below, the one-dimensional objects of string theory find a possible physical home in the \emph{dark sector}---the $\sim 85\%$ of the universe's matter content that interacts gravitationally but not electromagnetically---while visible-sector particles are modeled as compact vortex knots, topological defects with aspect ratios of order unity, more naturally described by knot theory than by string theory.

We emphasize at the outset that what follows is a phenomenological framework, not a UV-complete derivation.
Several of its central constructions---the duality between the visible and dark sectors, the spectrum-generating mode-locking mechanism, and the candidate string-theoretic correspondences of \S\ref{sec:dictionary}---remain conjectural pending a more formal field-theoretic treatment.
We have tried throughout to keep the boundary between calculation, model assumption, and conjecture explicit.

This proposal rests on the Null Worldtube Theory (NWT), a framework in which the vacuum is a superfluid condensate and particles are quantized vortex defects~\cite{NWT1,NWT2,NWT3,NWT4,NWT5,NWT6,NWT7,NWT8,NWT9}.
Previous work in the NWT series derived the Standard Model gauge group SU(3)$\times$SU(2)$\times$U(1) from the crossing algebra of vortex knots~\cite{NWT8}, predicted 56 particle masses with median error $0.40\%$ from four quantum numbers~\cite{NWT6}, and obtained the fine-structure constant, Weinberg angle, CKM and PMNS mixing matrices, and coupling constants from a single geometric parameter $\kappa = \pi^2$~\cite{NWT9}.

In this paper, we extend the framework to the dark sector by introducing a second condensate field, $\psi_\text{GUT}$, the residual of the GUT-scale condensate that did not fully reconnect during symmetry breaking.
The two condensates interpenetrate like the superfluid and normal components of helium~II, coupled only through gravity and a \emph{portal interaction}---a small, mediator-suppressed coupling between the visible and dark sectors of the kind familiar from hidden-sector dark-matter models~\cite{BerezhianiKhoury2015}, here arising as the cross-term in the two-condensate energy functional and quantified in \S\ref{sec:two-condensate}.
The dark sector's particle spectrum, cosmic string population, and contribution to dark energy all emerge from the properties of $\psi_\text{GUT}$.

Three observational anchors motivate the construction:
\begin{enumerate}
\item The cosmic string candidate CSc-1~\cite{Sazhina2019,Bulygin2023,Sazhina2025}, with inferred tension $G\mu/c^2 \sim 10^{-7}$, consistent with a GUT-scale vortex line.
\item The $\sim 1{,}000$ radio filaments near Sgr~A*~\cite{YusefZadeh2022}, perpendicular to the Galactic plane and remarkably straight---consistent with quantized vortex lines stirred by the black hole's rotation.
\item The $20$~GeV gamma-ray halo excess reported by Totani~\cite{Totani2025} from 15 years of Fermi-LAT data, with morphology matching dark matter annihilation in an NFW halo.
\end{enumerate}

We argue that all three observations admit a unified interpretation within the two-condensate framework, alongside the Bullet Cluster's dark matter morphology, the null results of direct-detection experiments, and the structure of the string theory landscape itself.

The paper is organized as follows.
Section~\ref{sec:two-condensate} introduces the two-condensate vacuum and derives the interface coupling.
Section~\ref{sec:cosmic-strings} identifies cosmic strings and galactic-center filaments as $\psi_\text{GUT}$ vortex lines.
Section~\ref{sec:dark-matter} derives the dark particle spectrum and confronts it with dark matter observations.
Section~\ref{sec:dictionary} presents the concept dictionary mapping string theory to NWT.
Section~\ref{sec:gauge} summarizes the gauge-group derivation and GUT extension.
Section~\ref{sec:veneziano} states a conjecture connecting the Veneziano amplitude to vortex reconnection.
Section~\ref{sec:dark-energy} interprets recent constraints on cosmic-string dark energy.
Section~\ref{sec:predictions} collects the predictions, and Section~\ref{sec:discussion} discusses implications and future directions.

%% ================================================================
\section{The two-condensate vacuum}
\label{sec:two-condensate}
%% ================================================================

\subsection{The Standard Model condensate $\psi_\text{SM}$}

In the NWT framework, the vacuum is modeled as a superfluid described by a complex scalar order parameter $\psi_\text{SM}$ satisfying the Gross--Pitaevskii equation (GPE):
\begin{equation}
i\hbar\frac{\partial\psi}{\partial t} = -\frac{\hbar^2}{2m^*}\nabla^2\psi + g|\psi|^2\psi,
\label{eq:GPE}
\end{equation}
where $m^* = m_e/\sqrt{2}$ is the effective mass of the condensate quasiparticle and $g$ is the self-interaction coupling.
We must be explicit about the status of this equation in what follows.
The GPE is used here as an \emph{analogue model} that captures the topological content of vortex defects---their classification by winding numbers and knot invariants---and not as a fundamental theory of the vacuum.
A fully relativistic treatment would require a Lorentz-invariant scalar field theory (\emph{e.g.}, a relativistic Gross--Pitaevskii or nonlinear Klein--Gordon equation), and constructing one is a priority for future work; the present paper does not assume such a theory has been built.
Throughout, where we invoke relativistic properties of the underlying vacuum---in particular the equality $c_s = c$ that justifies the Landau criterion in \S\ref{sec:cosmic-strings} and \S\ref{sec:dark-matter}---we are appealing to the (yet-to-be-constructed) relativistic completion, not to the non-relativistic GPE sound speed.
This separation between the analogue model used for topology and the relativistic theory invoked for kinematics is the central methodological caveat of the paper.
The condensate supports topological defects---quantized vortex lines---whose configurations are classified by torus knot invariants.

The single free parameter is the aspect ratio of the vortex torus, $\kappa_\text{SM} = R/r = \pi^2 \approx 9.87$, from which the fine-structure constant follows~\cite{NWT9}:
\begin{equation}
\alpha = \frac{1}{\sqrt{2}\,\kappa_\text{SM}^2} = \frac{1}{\sqrt{2}\,\pi^4} \approx \frac{1}{137.8}.
\label{eq:alpha}
\end{equation}
The healing length $\xi_\text{SM} = \hbar/(m_e c) = 386$~fm sets the characteristic size of visible-sector particles.
The gauge group SU(3)$\times$SU(2)$\times$U(1) arises from the crossing exchange algebra of the trefoil knot (3 crossings $\to$ $\mathfrak{su}(3)$) and the Hopf link (2 crossings $\to$ $\mathfrak{su}(2) \oplus \mathfrak{u}(1)$)~\cite{NWT8}.

\subsection{The GUT condensate $\psi_\text{GUT}$}

We model the GUT-era vacuum ($E_\text{GUT} \sim 10^{16}$~GeV) as a single condensate $\psi_\text{SU(5)}$ whose vortex topology is the cinquefoil knot $5_1$, with crossing algebra generating $\mathfrak{su}(5)$ via the rank-tower extension of~\cite{NWT8}.
GUT symmetry breaking corresponds to the topological surgery
\begin{equation}
5_1 \;\longrightarrow\; 3_1 + \text{Hopf link},
\label{eq:GUT-surgery}
\end{equation}
where the trefoil $3_1$ carries SU(3) and the Hopf link carries SU(2)$\times$U(1).
This surgery is a vortex reconnection event in the cooling condensate.

The reconnection need not be complete.
Regions of the original $\psi_\text{SU(5)}$ that failed to reconnect persist as a \emph{residual GUT condensate}, which we denote $\psi_\text{GUT}$.
The post-surgery vacuum is therefore a two-component superfluid:
\begin{equation}
\psi_\text{vac} = \psi_\text{SM} \oplus \psi_\text{GUT},
\label{eq:two-condensate}
\end{equation}
analogous to the two-fluid model of superfluid $^3$He--$^4$He mixtures~\cite{Vollhardt1990}, in which two interpenetrating condensates with independent order parameters coexist in the same volume.

The GUT condensate has its own aspect ratio $\kappa_\text{GUT}$, related to the GUT coupling $\alpha_\text{GUT} = 1/40$ by the same geometric relation as Eq.~\eqref{eq:alpha}:
\begin{equation}
\alpha_\text{GUT} = \frac{1}{\sqrt{2}\,\kappa_\text{GUT}^2} \quad\Rightarrow\quad \kappa_\text{GUT} = \frac{1}{\sqrt{\sqrt{2}\,\alpha_\text{GUT}}} \approx 5.32.
\label{eq:kappa-GUT}
\end{equation}
The healing length is $\xi_\text{GUT} = \hbar c / E_\text{GUT} \approx 1.5 \times 10^{-32}$~m---roughly $10^3$ times the Planck length.
The origin of the 19-order-of-magnitude separation between $\xi_\text{SM}$ and $\xi_\text{GUT}$ is not addressed in this work; we treat $\kappa_\text{GUT}$ phenomenologically rather than deriving it from a more fundamental theory, and the question of why the two condensates exist with such disparate scales remains open.

\subsection{The interface coupling}

The energy functional for the two-condensate system includes a cross-term:
\begin{equation}
\mathcal{E}_\text{int} = g_{12}\int |\psi_\text{SM}|^2\,|\psi_\text{GUT}|^2\,d^3x,
\label{eq:cross-term}
\end{equation}
where $g_{12}$ is the inter-condensate coupling.
The associated dimensionless coupling constant is
\begin{equation}
\alpha_\text{dark} = \frac{1}{\sqrt{2}\,\kappa_\text{SM}\,\kappa_\text{GUT}} = \frac{1}{\sqrt{2} \times \pi^2 \times 5.32} \approx \frac{1}{74.2},
\label{eq:alpha-dark}
\end{equation}
which matches the geometric mean $\sqrt{\alpha_\text{SM}\,\alpha_\text{GUT}} = 1/74.0$ to $0.3\%$.
This numerical agreement should be regarded as an internal consistency check rather than as a prediction: $\alpha_\text{GUT}$ was input via Eq.~\eqref{eq:kappa-GUT} and the coincidence is what an exact $\mathbb{Z}_2$ symmetry between the two condensates would imply.
The same caveat applies to the other numerical agreements quoted in this paper, which we flag explicitly where they arise.
This coupling governs all portal interactions between the visible and dark sectors.

The physical interpretation is transparent: $\alpha_\text{SM}$ is the self-coupling of $\psi_\text{SM}$ (proportional to $\kappa_\text{SM}^{-2}$), $\alpha_\text{GUT}$ is the self-coupling of $\psi_\text{GUT}$ (proportional to $\kappa_\text{GUT}^{-2}$), and $\alpha_\text{dark}$ is the cross-coupling (proportional to $(\kappa_\text{SM}\,\kappa_\text{GUT})^{-1}$).
The interface coupling is fixed by the model rather than introduced as a free fit parameter.

\subsection{Why dark matter is dark}

The two condensates have healing lengths separated by 19 orders of magnitude: $\xi_\text{SM} = 386$~fm versus $\xi_\text{GUT} = 1.5 \times 10^{-32}$~m.
To produce a dark particle (nucleate a vortex ring in $\psi_\text{GUT}$), one must disturb the condensate at the scale $\xi_\text{GUT}$.
No Standard Model process can resolve this scale: even the Higgs boson, at $125$~GeV, probes length scales of $\sim 10^{-18}$~m---fourteen orders of magnitude too coarse.

This \emph{structural inaccessibility} offers a possible common explanation for three null results:
\begin{itemize}
\item The LHC has not produced dark matter because collider energies cannot resolve $\xi_\text{GUT}$.
\item Direct-detection experiments (LZ, XENONnT) see nothing because nuclear recoils probe femtometer scales, not $10^{-32}$~m.
\item The barrier is not weak coupling---$\alpha_\text{dark} = 1/74$ is not small---but \emph{scale mismatch} between the two condensates.
\end{itemize}

The only channels through which dark matter can be observed are gravitational (lensing, rotation curves, the Bullet Cluster) and through annihilation products, where dark particles interact with \emph{each other} and the released energy spills into $\psi_\text{SM}$ via the portal.

%% ================================================================
\section{Cosmic strings as $\psi_\text{GUT}$ vortex lines}
\label{sec:cosmic-strings}
%% ================================================================

A quantized vortex line in $\psi_\text{GUT}$ is a one-dimensional topological defect with core radius $\xi_\text{GUT} \approx 1.5 \times 10^{-32}$~m and energy per unit length (tension)
\begin{equation}
\mu = \pi\rho_0\,\xi_\text{GUT}^2\,c^2\,\ln(\kappa_\text{GUT}),
\label{eq:string-tension}
\end{equation}
where $\rho_0$ is the condensate density.
In dimensionless form, the tension is related to the GUT energy scale via $G\mu/c^2 = (E_\text{GUT}/E_\text{Pl})^2$, giving
\begin{equation}
G\mu/c^2 \sim \alpha_\text{GUT} \sim 10^{-7}\text{--}10^{-6}
\end{equation}
for $E_\text{GUT} \sim 10^{15.9}\text{--}10^{16}$~GeV.
Within the present model, this identifies the corresponding defect as a \emph{cosmic-string candidate}---an object predicted independently by Kibble~\cite{Kibble1976}, studied extensively in the context of GUT phase transitions~\cite{Vilenkin1985}, and shown to arise naturally in superstring compactifications~\cite{Witten1985,Polchinski2005}.

An important distinction must be drawn between the tension of \emph{individual strings} and the constraints on a \emph{string network}.
The most stringent CMB bound, from combined Planck and ACT DR6 data, is $G\mu < 3.66 \times 10^{-8}$ at $2\sigma$ for a Nambu--Goto scaling network~\cite{Planck2013strings,ACT-DR6-strings}.
This bound is an order of magnitude below the individual CSc-1 tension.
Rather than excluding the model outright, this motivates the distinction between an individual string and a decayed network: in our framework the $\psi_\text{GUT}$ string network does \emph{not} persist in scaling form.
Quantitatively, the Planck+ACT bound is satisfied provided the surviving line-density obeys
\begin{equation}
f_\text{survive} \;\lesssim\; \frac{G\mu_\text{bound}}{G\mu_\text{individual}}
\;\sim\; \frac{3.66 \times 10^{-8}}{5 \times 10^{-7}} \;\approx\; 0.07
\label{eq:fsurvive}
\end{equation}
relative to a standard Nambu--Goto scaling configuration.
This is easily achieved.
Kelvin instability shreds a string on the microscopic timescale $\tau_K \sim \xi_\text{GUT}/c \sim 5 \times 10^{-41}$~s once thermal helical amplitudes reach $A/\lambda \sim 1$, which occurs essentially immediately after the GUT phase transition.
Compared with the cosmic time at recombination ($t_\text{rec} \sim 10^{13}$~s), the timescale ratio $t_\text{rec}/\tau_K \sim 10^{53}$ guarantees that any segment that ever entered the Kelvin-unstable regime has been converted to vortex rings (dark matter) many orders of magnitude before nucleosynthesis.
What survives at recombination is a sparse population of long-wavelength flat segments whose helical excitation was suppressed by random initial conditions, with collective energy density well below the CMB bound ($\Omega_s \sim 10^{-5}$ at $z = 0$; see \S\ref{sec:dark-energy}).
Similarly, PTA constraints ($G\mu \sim 10^{-10}\text{--}10^{-11}$ from the NANOGrav stochastic background~\cite{NANOGrav2023,Ellis2021}) apply to the gravitational wave emission of a scaling network, not to individual surviving strings.
The CSc-1 candidate, if confirmed, would be one such rare survivor---a single string whose individual tension $G\mu \sim 10^{-7}$ is consistent with the observed lensing signal, while the network as a whole has decayed below the CMB and PTA thresholds.

\subsection{The CSc-1 candidate}

The most developed cosmic string candidate is CSc-1, identified from CMB anisotropy analysis and studied through optical searches for gravitational lens chains~\cite{Sazhina2019}.
Deep photometry of the double galaxy SDSS~J110429.61+233150.3 in the CSc-1 field reveals spectral consistency between the two components, suggestive of gravitational lensing by a linear mass distribution.

Bulygin, Sazhin \& Sazhina~\cite{Bulygin2023} developed the complete lensing theory for a curved cosmic string, demonstrating that string inclination and bending produce asymmetric image pairs with position-angle offsets---a signature absent for straight strings.
Their simulations employ $G\mu = 7.0 \times 10^{-7}$, yielding a critical bending angle $\theta_c \approx 13^\circ$ beyond which the second image vanishes.
This explains the observational deficit of lens chains: most string segments have bends exceeding $\theta_c$.

NWT identifies CSc-1 as a vortex line in $\psi_\text{GUT}$ with tension $G\mu/c^2 \in [4, 7] \times 10^{-7}$, corresponding to $E_\text{GUT} \in [7.7, 10.2] \times 10^{15}$~GeV---squarely within the standard GUT window.
The string is curved by the gravitational potential through which it passes, producing the asymmetric lensing predicted by Bulygin et al.
The most recent review of observational methods~\cite{Sazhina2025} confirms that CSc-1 remains the leading candidate and that new search strategies based on curved-string theory significantly expand the discoverable parameter space.

\subsection{Galactic center filaments}

The MeerKAT radio telescope has revealed $\sim 1{,}000$ filamentary structures within $\sim 150$~ly of Sgr~A*~\cite{YusefZadeh2022}.
The \emph{vertical} filaments are perpendicular to the Galactic plane, remarkably straight (up to $\sim 50$~pc in length), and show amplified magnetic fields along their axes.
They appear in pairs and clusters with regular spacing of $\sim 1$~AU within clusters.

In NWT, these are quantized vortex lines in $\psi_\text{GUT}$, created by Sgr~A*'s rotation---the astrophysical analogue of Feynman's rotating-bucket experiment in superfluid helium.
Three qualitative predictions match observations:
\begin{enumerate}
\item \textbf{Perpendicularity.} Vortex lines in a rotating superfluid align with the rotation axis (here, Sgr~A*'s spin axis $\approx$ Galactic north).
\item \textbf{Straightness.} The condensate density ratio $\rho_\text{GUT}/\rho_\text{ISM} \sim 10^{24}$ suppresses Kelvin--Helmholtz instability; the vortex lines are effectively rigid rods that nothing in the Galactic center environment can bend.
\item \textbf{Synchrotron emission.} The interstellar medium acts as a tracer: plasma trapped along the vortex lines emits synchrotron radiation, making the otherwise invisible $\psi_\text{GUT}$ defects observable in radio.
\end{enumerate}

The $\sim 1$~AU spacing between paired filaments admits a striking explanation.
Since vortex lines cannot terminate in the bulk of a superfluid (the winding number is a topological invariant), each filament must form a closed loop.
The natural geometry is a \emph{hairpin}: a vortex ring nucleated at the ergoregion (diameter $\sim R_s$) is expelled along the spin axis by the Magnus force in the rotating condensate, and Keplerian shear stretches it vertically into an elongated loop whose two legs appear as a pair of parallel filaments.
The hairpin width---the separation between the two legs---is set by the diameter of the ergoregion at nucleation:
\begin{equation}
d_\text{pair} \approx R_s(\text{Sgr~A*}) = \frac{2GM}{c^2} \approx 0.08\text{~AU}.
\end{equation}
With modest expansion ($\sim 13\times$) as the ring rises through the decreasing-density environment, $d_\text{pair} \sim 1$~AU---matching the observed spacing.
The loop closes smoothly through the ergoregion below the accretion disk, with frame-dragging continuously driving circulation through the loop.

This model makes a sharp scaling prediction: $d_\text{pair} \propto M_\text{BH}$.
For M87* ($M = 6.5 \times 10^9\;M_\odot$), the predicted pair spacing is $\sim 1{,}600\times$ larger than Sgr~A*'s, or $\sim 0.1$~pc---potentially resolvable with VLBI.

Dark matter production occurs at the \emph{hairpin bend}---the top of the loop, where curvature is maximum and Kelvin instability is strongest.
Each bend continuously sheds vortex rings (dark particles), making the filaments dark matter \emph{fountains}: Sgr~A*'s rotation $\to$ vortex loop creation $\to$ outward propagation $\to$ ring shedding at the tips.
JWST observations reveal that Sgr~A*'s emission fluctuates continuously at timescales from seconds to months~\cite{YusefZadeh2025JWST}, consistent with ongoing stochastic vortex nucleation at the horizon.

The stability of the filaments, compared to primordial strings that shredded in the early universe, follows from the thermal environment: at $kT/E_\text{vortex} \sim 10^{-31}$, thermal Kelvin-wave excitation is negligible.
Primordial strings at $T \sim 10^{28}$~K had $A/\lambda \sim 10^6$, driving immediate self-reconnection into vortex rings (dark matter particles).

The population study by Yusef-Zadeh et al.~\cite{YusefZadeh2022spacing} reveals that grouped filaments show ``harp-like, fragmented cometary tail-like, or loop-like structures'' with ``striking examples of a single filament splitting into two prongs at a junction''---precisely the hairpin morphology.
The bimodal position angle distribution~\cite{YusefZadeh2023horizontal}---vertical filaments (PA $\sim 0^\circ$, nonthermal synchrotron) and horizontal filaments (PA $\sim 70^\circ$, thermal, radial toward Sgr~A*)---indicates two mechanisms: the vertical population from steady-state spin-driven hairpin loops, and the horizontal population from a transient collimated outflow $\sim 6$~Myr ago with opening angle $\sim 40^\circ$.

Swamy \& Singh~\cite{Swamy2025} have shown that a cosmic string attached to a rotating black hole extracts angular momentum via the torque $\dot{H}_\text{CS} = -(\mu/M)[\mathbf{J} - (\hat{\mathbf{j}}\cdot\mathbf{J})\hat{\mathbf{j}}]$, providing a quantitative framework for the BH--string interaction that sustains the hairpin loops.
Their broader program~\cite{Swamy2025accretion,Swamy2025QNM,Swamy2025Xray,Swamy2025DM} establishes that BH--string systems produce observable effects in jet power, quasinormal mode spectra, X-ray binary timing, and dark matter genesis---all consistent with the NWT hairpin picture.

\subsection{String stability: why GUT strings survive}

Not all vortex strings persist to the present day.
QCD-scale strings ($E \sim 170$~MeV) undergo Kelvin instability at the hadronization temperature: the ratio of thermal amplitude to string wavelength exceeds unity ($A/\lambda \sim 10^6$), and the strings self-reconnect into vortex rings (hadrons) on timescales of $\sim 10^{-23}$~s.
This \emph{is} hadronization in the NWT picture: the QGP-to-hadron transition is vortex filament shredding.

GUT-scale strings, by contrast, have $A/\lambda \ll 1$ at all post-inflationary temperatures.
The Landau critical velocity for $\psi_\text{GUT}$ is $v_\text{Landau} = c_s = c$ (the condensate is Lorentz-invariant), and all astrophysical motions are deeply subsonic.
GUT strings are stable against Kelvin instability, thermal buffeting, and hydrodynamic drag at late times.
However, the primordial network decays through three complementary mechanisms:
(i)~\emph{Kelvin instability} (the NWT mechanism): thermal excitation of helical oscillations $\to$ self-reconnection $\to$ vortex ring shedding;
(ii)~\emph{monopole-antimonopole nucleation}~\cite{Tranchedone2026}: metastable strings break by nucleating topological defect pairs, with the rate exponentially enhanced at finite temperature---Tranchedone et al.\ show that ``metastable cosmic strings are broken at the start'';
(iii)~\emph{semiclassical evaporation}~\cite{Penna2025}: a Hawking-like process producing particles from the string core.
All three convert string energy into particles, gradually transforming the network into dark matter (vortex rings) and gravitational radiation.
The fragmentation can cascade: loops produced by the network break into smaller loops~\cite{Auclair2025}, which break further, until the smallest topologically stable loops---at the Compton scale of the dark particle spectrum---persist as dark matter.


%% ================================================================
\section{Dark matter as closed strings}
\label{sec:dark-matter}
%% ================================================================

A vortex ring in $\psi_\text{GUT}$ is a closed loop of vortex line---naturally described within this framework as a closed-string-like defect.
Its core has width $\xi_\text{GUT} \approx 10^{-32}$~m, while its ring radius is set by the particle's Compton wavelength, giving aspect ratios $R/\xi \sim 10^{13}\text{--}10^{17}$---thin loops that are effectively one-dimensional at all scales above the GUT healing length.

\subsection{The dark particle spectrum}

The $\psi_\text{GUT}$ condensate supports the same mode-locking spectrum as $\psi_\text{SM}$~\cite{NWT6}, but with $\kappa_\text{GUT} = 5.32$ replacing $\kappa_\text{SM} = \pi^2$.
Setting the dark proton mass to $m_{d\text{-}p} = 187$~GeV (from the relic abundance condition with $\sigma = \alpha^3\pi\lambdabar^2$~\cite{NWT4}), the mass ratio $m_{d\text{-}p}/m_{d\text{-}e} = 2020$ from the mode-locking formula determines the full dark spectrum:

\begin{table}[h]
\centering
\caption{Dark particle spectrum from $\psi_\text{GUT}$ mode-locking ($\kappa_\text{GUT} = 5.32$).  Stability assumes the dark weak force is broken at the GUT scale.}
\label{tab:dark-spectrum}
\begin{tabular}{@{}llccc@{}}
\toprule
Symbol & Name & Mass & Stable? & Role \\
\midrule
$d$-$e$    & dark electron & 93~MeV   & \checkmark & lightest dark lepton \\
$d$-$\nu$  & dark neutrino & $\sim 0$ & \checkmark & phase soliton \\
$d$-$\mu$  & dark muon     & 18.8~GeV & \checkmark & Totani signal \\
$d$-$\pi$  & dark pion     & 25~GeV   & \checkmark & dark nuclear force \\
$d$-$K$    & dark kaon     & 93~GeV   & $\times$   & $\to d\text{-}\pi$ \\
$d$-$p$    & dark proton   & 187~GeV  & \checkmark & dominant $\Omega_\text{DM}$ \\
$d$-$n$    & dark neutron  & 187~GeV  & \checkmark & isospin partner \\
\bottomrule
\end{tabular}
\end{table}

The relic abundance scales as $\Omega_i \propto m_i^3$ for geometric cross-sections, so the dark proton dominates at $99.8\%$ of $\Omega_\text{DM}$.
The lighter species ($d$-$e$, $d$-$\mu$, $d$-$\pi$) are cosmologically subdominant but produce distinctive annihilation signatures.

\subsection{The Totani 20~GeV excess}

Totani~\cite{Totani2025} reports a statistically significant excess in the Fermi-LAT gamma-ray data, peaking at $\sim 20$~GeV with halo-like (NFW-$\rho^2$) morphology consistent with dark matter annihilation.
The signal is detected at $|b| \geq 10^\circ$ and is consistent across sky quadrants.
The inferred WIMP mass under a pure $b\bar{b}$ channel is $500$--$800$~GeV, with cross-section $(5\text{--}8) \times 10^{-25}$~cm$^3$/s.

The model allows two possible interpretations within the multi-species dark sector:
\begin{itemize}
\item \textbf{Dark proton continuum.} The $d$-$p$ at 187~GeV annihilating via the Higgs portal into a mix of $\sim 31\%$ $b\bar{b}$ + $69\%$ $WW$ produces a gamma-ray spectrum peaking at $\sim 20$~GeV.
\item \textbf{Dark muon line.} The $d$-$\mu$ at 18.8~GeV annihilating to $\gamma\gamma$ produces a monochromatic line at 18.8~GeV---within $6\%$ of the Totani peak.  Fermi-LAT's $\sim 10\%$ energy resolution at 20~GeV would smear this into a narrow excess indistinguishable from the observed spectrum.
\end{itemize}

The cross-section enhancement over the canonical thermal relic value ($2.5 \times 10^{-26}$~cm$^3$/s) is provided by Sommerfeld enhancement with $\alpha_\text{dark} = 1/74$:
\begin{equation}
S = \frac{\pi\alpha_\text{dark}}{v_\text{halo}/c} \approx \frac{\pi/74}{7.3 \times 10^{-4}} \approx 58,
\end{equation}
giving $\langle\sigma v\rangle_\text{eff} \sim 1.4 \times 10^{-24}$~cm$^3$/s, within a factor of two of the Totani value after accounting for the dark muon's subdominant number density.

An important constraint comes from AMS-02 antiproton data.
If the $d$-$\mu$ annihilates primarily to $b\bar{b}$, the resulting antiproton flux would be in tension with AMS-02 bounds at $\sim 20$~GeV.
The NWT topology suggests an escape: the $d$-$\mu$ is a \emph{dark lepton} (genus $g = 0$, $n_q = 0$), and a leptophilic annihilation channel $d$-$\mu\,\overline{d\text{-}\mu} \to \gamma\gamma$ would produce photons but not antiprotons, with the hadronic channels $(b\bar{b}, WW)$ kinematically open but topologically suppressed.
We flag this as a critical theoretical uncertainty, not a resolved point: the full portal matrix element remains to be computed, and the dark-muon interpretation of the Totani signal stands or falls on its outcome.
Until that calculation is performed, the dark-proton continuum interpretation should be regarded as the conservative reading of the same data.

\subsection{The Bullet Cluster as a subsonic superfluid collision}

The Bullet Cluster (1E~0657-56) is the textbook case for collisionless dark matter: the DM component, traced by gravitational lensing, passes through the collision while the baryonic gas shocks and lags behind~\cite{JWST-Bullet2025}.
The standard interpretation is that DM particles have negligibly small self-interaction cross-sections ($\sigma/m < 1$~cm$^2$/g).

NWT provides a different explanation with the same observational consequences.
The collision velocity of $\sim 4{,}500$~km/s corresponds to Mach number
\begin{equation}
\mathcal{M} = \frac{v}{c_s} = \frac{4{,}500\text{ km/s}}{c} \approx 0.015
\end{equation}
in the $\psi_\text{GUT}$ condensate, where $c_s = c$ (the condensate is Lorentz-invariant).
This is \emph{deeply subsonic}.
The Landau criterion for superfluidity---no dissipation for flow speeds below $c_s$---is satisfied by a factor of $\sim 67$.
The two $\psi_\text{GUT}$ halos simply pass through each other without shock, drag, or vortex nucleation.

``Collisionless dark matter'' is, in NWT, ``subsonic superfluid dark matter.''
The physics is identical in its observational consequences at cluster scales---we do not claim a distinction that current data could resolve.
The difference is ontological and mechanistic: the self-interaction cross-section need not be small; the sound speed need only be large.
The testable consequences of the superfluid picture emerge at other scales---through cosmic string lensing (\S\ref{sec:cosmic-strings}), dark particle annihilation (\S\ref{sec:dark-matter}), and the string network's contribution to dark energy (\S\ref{sec:dark-energy})---not at the cluster scale itself.
We note that the superfluid dark matter literature~\cite{BerezhianiKhoury2015} typically invokes phonon-mediated forces to produce MOND-like behavior in galaxies; the NWT condensate, with $c_s = c$, does not produce such forces and is indistinguishable from CDM at all astrophysical scales above $\xi_\text{GUT}$.

JWST observations~\cite{JWST-Bullet2025} reveal additional morphological features consistent with the superfluid picture:
\begin{itemize}
\item A \emph{bridge} of elevated DM density between the two mass peaks (the overlap region of two interpenetrating halos).
\item Smooth, shock-free DM gradients (no discontinuities, unlike the gas).
\item DM--BCG offsets of 44--70~kpc, consistent with a dynamical time lag $\sim R_\text{virial}/v$.
\end{itemize}

\subsection{Direct detection and self-interaction}

The structural inaccessibility discussed in \S\ref{sec:two-condensate} implies that the spin-independent nuclear scattering cross-section is suppressed far below current experimental bounds.
The dark proton's self-interaction cross-section is $\sigma/m \sim \pi\lambdabar_{d\text{-}p}^2/m_{d\text{-}p} \sim 10^{-10}$~cm$^2$/g, safely below the Bullet Cluster constraint by ten orders of magnitude.

%% ================================================================
\section{The concept dictionary: string theory $\leftrightarrow$ NWT}
\label{sec:dictionary}
%% ================================================================

The mathematical structures of string theory and NWT are closely related, and we explore them as candidate descriptions of a common underlying vortex physics---vortex defects in a superfluid vacuum---viewed in two different aspect-ratio limits.
Table~\ref{tab:dictionary} presents the proposed correspondence.

\begin{table}[h]
\centering
\caption{Concept dictionary mapping string theory to NWT.  The two frameworks describe the same vortex physics in different aspect-ratio regimes: $R/\xi \sim 10^{14}$ (strings, dark sector) versus $R/\xi \sim 1$ (knots, visible sector).}
\label{tab:dictionary}
\begin{tabular}{@{}lll@{}}
\toprule
\textbf{String theory} & \textbf{NWT} & \textbf{Physical content} \\
\midrule
Fundamental string & Vortex line in $\psi_\text{GUT}$ & 1D topological defect \\
String tension $\alpha'$ & Vortex tension $\mu$ & Energy per unit length \\
String length $l_s = \sqrt{\alpha'}$ & Healing length $\xi_\text{GUT}$ & Characteristic scale \\
Closed string & Vortex ring (dark particle) & Closed 1D loop \\
Open string & Vortex line with endpoints & Ends on condensate boundary \\
Worldsheet & Vortex tube surface & 2D surface swept by string \\
Worldsheet CFT & Helmholtz eigenmodes on tube & 2D field theory on surface \\
\midrule
Calabi-Yau 3-fold & Torus $T^2$ with aspect ratio $\kappa$ & Compactification manifold (analogue) \\
CY moduli & $\kappa_\text{SM} = \pi^2$, $\kappa_\text{GUT} = 5.32$ & Shape parameters \\
Moduli stabilization & Phase closure ($m \in \mathbb{Z}$) & Discrete spectrum from geometry \\
Kaluza-Klein tower & $m = 3, 4, 5, \ldots$ excitations & Tower of massive states \\
\midrule
Gauge group from CY topology & Gauge group from crossing algebra & $N$ crossings $\to$ SU($N$) \\
Chan-Paton factors & Strand labels at crossings & Internal symmetry indices \\
D-brane & Condensate sheet ($\psi_\text{SM}$ or $\psi_\text{GUT}$) & Extended object hosting open strings \\
Brane separation & Higgs VEV $v_h = 246$~GeV & Inter-sector coupling scale \\
\midrule
$10^{500}$ landscape & Countable $(p,q,m,n_q,f)$ modes & Vacuum multiplicity \\
Supersymmetry & SM $\leftrightarrow$ dark correspondence & Partner spectrum (ratio $\sim 200$) \\
\bottomrule
\end{tabular}
\end{table}

\subsection{Six extra dimensions as six quantum numbers}

String theory requires ten spacetime dimensions: 3+1 observed plus 6 compactified on a Calabi-Yau manifold.
In NWT, the six compactified dimensions find a suggestive analogue in six discrete quantum numbers that label the torus knot state:
\begin{itemize}
\item $p$ (toroidal winding) and $q$ (poloidal winding): two angular coordinates on $T^2$.
\item $m$ (phase-closure integer): the Kaluza-Klein excitation level.
\item $n_q$ (constituent count from carrier crossing number): determines the confinement structure.
\item $f$ (framing): encodes electric charge via $Q = f/2 + (B+S)/2$.
\item $g$ (genus): distinguishes leptons ($g = 0$) from hadrons ($g > 0$).
\end{itemize}

We do not claim these are literal identifications---the discrete quantum numbers lack the continuous geometric structure (metric, holonomy, local propagation) of true compactified dimensions.
Rather, each quantum number labels an independent degree of freedom that, in string theory, would correspond to a winding or wrapping mode on a cycle of the internal manifold.
The torus knot parameter space plays a role \emph{analogous} to the Calabi-Yau moduli space, with phase closure ($m \in \mathbb{Z}$) serving as an NWT analogue of flux stabilization.
Whether this analogy can be sharpened into a precise mathematical correspondence remains an open question.

The dramatic simplification relative to standard string compactification is evident: a typical Calabi-Yau has $\mathcal{O}(100)$ moduli giving rise to the $10^{500}$ landscape.
NWT has \emph{one} continuous parameter ($\kappa$) and five discrete quantum numbers, yielding a countable spectrum---a periodic table rather than a landscape.

\subsection{The string/knot duality}

The deepest entry in the dictionary is the duality between strings and knots, which distinguishes the dark and visible sectors:

\begin{center}
\begin{tabular}{@{}lcc@{}}
\toprule
& \textbf{Visible ($\psi_\text{SM}$)} & \textbf{Dark ($\psi_\text{GUT}$)} \\
\midrule
Topology & Compact knots & Thin loops (strings) \\
Aspect ratio $R/\xi$ & $\sim 1$ & $\sim 10^{14}$ \\
Mathematics & Knot theory & String theory \\
\bottomrule
\end{tabular}
\end{center}

We propose this as a \emph{conjectured correspondence}---not yet a duality in the rigorous sense (which would require a map between Hilbert spaces, preservation of correlation functions, and a controlled limit), but a candidate relationship between two aspect-ratio regimes of the same vortex physics.
Establishing whether this correspondence can be promoted to a genuine duality---with, for instance, a partition function equivalence $Z_\text{knot}[\text{data}] = Z_\text{string}[\text{mapped data}]$---is a central open problem.
If it holds, string theory and knot theory would not be competing descriptions of nature but \emph{two limits} of vortex defect topology, united by the superfluid vacuum.

\subsection{Candidate string-theoretic dualities}
\label{sec:candidate-dualities}

A central feature of string theory is its web of dualities: T-duality, S-duality, and AdS/CFT, each relating apparently distinct string configurations as physically equivalent descriptions of the same underlying theory.
NWT has structural analogues of each.
We list them here as targets for future analysis rather than as established results, because in each case the analogy is at the level of the parameter exchange, not at the level of equivalence of correlation functions or partition functions.

\paragraph{T-duality.}
String-theoretic T-duality exchanges a compactification radius $R$ with $\alpha'/R$, mapping winding modes to momentum modes on the dual circle.
In NWT, the structural analogue is the exchange $\kappa_\text{SM} \leftrightarrow \kappa_\text{GUT}$: the small-aspect-ratio limit (compact knots, visible sector, $R/\xi \sim 1$) is mapped to the large-aspect-ratio limit (thin loops, dark sector, $R/\xi \sim 10^{14}$), with the toroidal and poloidal winding numbers $(p, q)$ playing roles formally analogous to momentum and winding modes.
The numerical content $\kappa_\text{SM} \cdot \kappa_\text{GUT} \approx \pi^2 \times 5.32 \approx 52.5$ is suggestive of a fixed self-dual scale, but no proof of mode-spectrum equivalence has been constructed and the analogy should not be over-read.

\paragraph{S-duality.}
S-duality exchanges weak and strong coupling, $g_s \leftrightarrow 1/g_s$, mapping a perturbatively tractable theory to its non-perturbative dual.
In NWT, the structural analogue is the exchange $\alpha_\text{SM} \leftrightarrow \alpha_\text{GUT}$: the visible-sector coupling at $\alpha_\text{SM} \approx 1/137$ is mapped to the dark-sector coupling at $\alpha_\text{GUT} \approx 1/40$.
The interface coupling $\alpha_\text{dark} = 1/(\sqrt{2}\,\kappa_\text{SM}\kappa_\text{GUT})$ matches the geometric mean $\sqrt{\alpha_\text{SM}\,\alpha_\text{GUT}}$ to $0.3\%$ (Eq.~\eqref{eq:alpha-dark}), the precise relation that an exact $\mathbb{Z}_2$ S-duality acting on the two-condensate vacuum would imply, with the interface coupling sitting at the self-dual midpoint.
We do not claim this constitutes a proof; the numerical match is the principal evidence and the underlying invariance has not been demonstrated.

\paragraph{AdS/CFT.}
The AdS/CFT correspondence relates a gravitational theory in the bulk of an asymptotically anti-de Sitter spacetime to a conformal gauge theory on its boundary.
In NWT, the structural analogue is the bulk--boundary relation between the two condensates: $\psi_\text{GUT}$ (strings, gravitationally relevant by virtue of GUT-scale tension) plays the role of the bulk theory, while $\psi_\text{SM}$ (knots, gauge-theoretic by virtue of the crossing algebra of \S\ref{sec:gauge}) plays the role of the boundary theory.
This is the most speculative entry of the three: the explicit form of the bulk metric, the boundary CFT action, and the matching of correlation functions all remain to be developed.
We mention the analogy because the structural picture---a gravitational string sector dual to a gauge-theoretic knot sector, related by aspect ratio---fits the AdS/CFT pattern unusually closely for an ansatz constructed from condensed-matter intuition.

These three correspondences extend the more concrete dictionary entries (gauge group from topology, mass spectrum from mode-locking, gauge coupling from $\kappa$) into territory where NWT does not yet have rigorous results.
They are listed to make the targets explicit, in the hope that practitioners trained in string-theoretic methods may sharpen them.


%% ================================================================
\section{Gauge groups from crossing algebra}
\label{sec:gauge}
%% ================================================================

In a previous paper~\cite{NWT8}, we argued that the Standard Model gauge group SU(3)$\times$SU(2)$\times$U(1) is the crossing exchange algebra of the two simplest vortex topologies: the trefoil knot (3 self-crossings $\to$ $\mathfrak{su}(3)$) and the Hopf link (2 inter-crossings $\to$ $\mathfrak{su}(2)$, linking phase $\to$ $\mathfrak{u}(1)$).
The mechanism is as follows: at each crossing, continuous strand exchange defines an operator analogous to a transposition in the symmetric group.
The $N$ crossing operators serve as \emph{seed generators}; the full $\mathfrak{su}(N)$ algebra (dimension $N^2 - 1$) emerges through successive commutators, not directly from the $N$ crossings.
For the trefoil ($N = 3$), three seed operators are identified explicitly in~\cite{NWT8} with the real off-diagonal Gell-Mann matrices $\lambda_1/2$, $\lambda_4/2$, $\lambda_6/2$, and all eight Gell-Mann matrices are then recovered through two levels of nested commutators.
The Hopf link case ($\mathfrak{su}(2) \oplus \mathfrak{u}(1)$) is worked out analogously.

We are careful here about what is and is not proved in~\cite{NWT8}.
The closure result for $N = 3$ is established by explicit construction.
The extension to higher $N$---in particular to the cinquefoil ($N = 5$) carrier of $\mathfrak{su}(5)$ invoked in the GUT-breaking discussion above---is not worked out crossing-by-crossing in~\cite{NWT8}.
Instead, it is invoked via the standard Lie-theoretic result that the real off-diagonal generators $\{T_{ij}\}_{i<j}$, viewed as elementary matrices acting on $\mathbb{C}^N$, generate all of $\mathfrak{su}(N)$ for any $N \geq 2$, together with the identification of the cinquefoil's five crossings with five such off-diagonal generators in the obvious basis.
The physical claim of the present paper is the \emph{identification} of these abstract generators with vortex-knot crossing exchanges; the algebra closure itself follows from textbook Lie theory and is independent of any topological interpretation.
We continue to regard the further questions---invariance under Reidemeister moves, the precise relation to braid group representations and quantum-group structures, and an explicit crossing-level closure construction for general $N$---as open problems requiring additional mathematical work, and we do not claim them as resolved here.

The crossing algebra extends naturally to the GUT scale.
Identifying the five crossings of the cinquefoil knot $5_1$ with five real off-diagonal generators of $\mathfrak{su}(5)$ in the obvious basis on $\mathbb{C}^5$, the resulting algebra closes on the full 24-dimensional $\mathfrak{su}(5)$ via the standard Lie-theoretic result invoked above.
The strand decomposition $\{0,1,2\} + \{3,4\}$ reproduces the Standard Model embedding SU(3)$\times$SU(2)$\times$U(1) $\subset$ SU(5), with hypercharge $Y = \text{diag}(1/3,\,1/3,\,1/3,\,-1/2,\,-1/2)$ emerging automatically from the $3+2$ partition---a textbook Georgi--Glashow result that the strand decomposition simply realizes geometrically.

The GUT breaking sequence~\eqref{eq:GUT-surgery} is then a knot simplification: the $5_1$ reconnects to $3_1$ + Hopf, shedding the 12 broken generators (the standard $24 = 8 + 3 + 1 + 12$ Georgi--Glashow decomposition) as massive $X$ and $Y$ bosons that we interpret as reconnection debris with mass $\sim E_\text{GUT}$.

An extension of this picture to the next GUT step,
\begin{equation}
\text{SO}(10) \;\to\; \text{SU}(5) \;\to\; \text{SU}(3)\times\text{SU}(2)\times\text{U}(1),
\end{equation}
seems natural in the present framework, although it has not been worked out in earlier NWT papers and is offered here as a conjectural extension rather than an established result.
We propose that the chain might correspond to $5_1 + 5_1^* \to 5_1 \to 3_1 + \text{Hopf}$, with SO(10) arising from the chirality doubling of the $5_1$ knot with its mirror image and the right-handed neutrino realized as the singlet in the $5_1^*$ sector.
A formal construction of the SO(10) crossing algebra, and a verification that the doubling reproduces the SO(10) representation content, remain open problems.

That the same knot invariants are now becoming experimentally accessible on quantum hardware lends an unexpected empirical anchor to the picture.
Following the proposal that any quantum computation can be reformulated as a Jones polynomial evaluation~\cite{Meichanetzidis2025}, we have computed Jones polynomials at the fifth root of unity for the six lightest NWT torus knots $T(p,q)$---the carrier topologies of the trefoil ($\pi^0$, $\eta$), $K^\pm$, $D^\pm$, $\Lambda$/$\tau$, $\pi^\pm$, and $D^0$ states---on IBM Heron-r2 superconducting hardware via a Temperley--Lieb path model implemented through the Hadamard test~\cite{NWTJonesQC}.
After dynamical decoupling and Pauli twirling, the six measured invariants reproduce the analytic values with errors of $3$--$18\%$ (median $\sim 8\%$), with the worst case being the $\pi^\pm$ knot whose invariant has the smallest magnitude and is therefore most sensitive to residual coherent error.
This is, to our knowledge, the first hardware computation of knot invariants for the topologies that classify visible-sector vortices in NWT, and provides a route to direct experimental discrimination among candidate carrier knots that does not rely on collider production.


%% ================================================================
\section{The Veneziano amplitude conjecture}
\label{sec:veneziano}
%% ================================================================

This section is offered as a research direction, not as a result.
No derivation of what follows is currently known to us; we record the conjecture because the geometric ingredients line up suggestively, and because identifying the most concrete target may help motivate the more formal work that would be required to settle it.

The Veneziano amplitude~\cite{Veneziano1968},
\begin{equation}
A(s,t) = \frac{\Gamma(-\alpha(s))\,\Gamma(-\alpha(t))}{\Gamma(-\alpha(s)-\alpha(t))},
\end{equation}
with Regge trajectory $\alpha(s) = \alpha_0 + \alpha' s$, was the first dual resonance model and led directly to the string interpretation of hadron scattering.

\textbf{Conjecture.} \emph{The Veneziano amplitude emerges from the path integral over vortex reconnection configurations in $\psi_\text{GUT}$.}

The supporting argument is geometric.
When two vortex rings approach, reconnect, and separate, the vortex lines sweep out a two-dimensional surface in spacetime---the \emph{worldsheet}.
The moduli space of reconnection configurations (parameterized by impact parameter, relative velocity, and winding-number matching) is a Riemann surface.
The path integral over this moduli space, weighted by the GPE action, should yield an amplitude with:
\begin{itemize}
\item Crossing symmetry ($s \leftrightarrow t$), inherited from the topological equivalence of reconnection channels.
\item Regge behavior ($J = \alpha' m^2$), from the rotational modes of the vortex ring.
\item Duality, from the equivalence between open and closed reconnection geometries.
\end{itemize}

A proof would establish a direct mathematical bridge between NWT vortex physics and the perturbative string S-matrix.
We leave this as an invitation to the string theory community.


%% ================================================================
\section{Cosmic strings as dynamical dark energy}
\label{sec:dark-energy}
%% ================================================================

Cheng, Di~Valentino \& Visinelli~\cite{Cheng2026} have recently constrained phenomenological models of a residual cosmic string network contributing to the late-time expansion, using Planck CMB data, DESI baryon acoustic oscillations, and Type~Ia supernovae.
Their Models~3 and 4, which allow the string density parameter $\Omega_s$ to take negative values, show a consistent preference for $\Omega_s < 0$---a negative effective energy density from the string network.

In NWT, this result admits a natural interpretation.
The $\psi_\text{GUT}$ vortex network \emph{screens} the cosmological constant: the string network's equation of state ($w_s = -1/3$ for frozen strings) partially cancels the vacuum energy.
As the network decays via Kelvin instability---shedding vortex rings (dark matter particles) over cosmic time---the screening diminishes and $\Lambda_\text{eff}$ increases.

This connects three cosmological phenomena through a single mechanism:
\begin{enumerate}
\item \textbf{Dark matter}: vortex rings shed from the decaying network.
\item \textbf{Dark energy}: the progressive loss of $\Lambda$ screening as the network thins.
\item \textbf{Cosmic strings}: the surviving network remnants (CSc-1).
\end{enumerate}

The coincidence problem---why $\Omega_\Lambda \sim \Omega_\text{matter}$ today---is resolved because the network decay timescale is set by $E_\text{GUT}$, which also determines the matter--radiation equality scale.
The model predicts $w_\text{eff} = -1 + \varepsilon(z)$ with $\varepsilon > 0$ decreasing over time as the last strings decay, consistent with the DESI hint of time-varying dark energy.

Using the one-scale model for the string network with loop-chopping efficiency $\varepsilon_\text{loop} \sim 50$, the present-day string density is $\Omega_s \sim 10^{-5}$---within the Cheng et al.\ bound of $|\Omega_s| < 0.009$ and consistent with CMB constraints at recombination ($\Omega_s(z\sim1100) \sim 0.01$, at the Planck limit).


%% ================================================================
\section{Predictions and observational tests}
\label{sec:predictions}
%% ================================================================

Table~\ref{tab:predictions} collects the principal predictions of the two-condensate model.

\begin{table}[h]
\centering
\caption{Quantitative predictions and their experimental tests.}
\label{tab:predictions}
\begin{tabular}{@{}p{3.2cm}p{2.5cm}p{2.0cm}@{}}
\toprule
\textbf{Prediction} & \textbf{Value} & \textbf{Test} \\
\midrule
Proton decay & $\tau_p \in [2.4, 5.5]\times 10^{34}$~yr$^{*}$ & Hyper-K \\
DM mass ($d$-$p$) & 187~GeV & LZ/XENONnT \\
Totani signal ($d$-$\mu$) & 18.8~GeV line & Fermi-LAT, CTA \\
$\gamma\gamma$ line & 18.8 or 187~GeV & CTA, AMEGO-X \\
String tension & $G\mu \in [4,7]\times 10^{-7}$ & CSc-1 spectroscopy \\
Bullet Cluster & Mach~0.015 & JWST lensing \\
GW background & $G\mu \sim 10^{-7}$ & LISA, ET, CE \\
AMS-02 $\bar{p}$ & $d$-$\mu$ at 19~GeV & AMS-02 \\
DM inaccessibility & No collider signal & LHC, FCC \\
\bottomrule
\end{tabular}\\[2pt]
\footnotesize $^{*}$ Lower bound set by the Super-Kamiokande exclusion $\tau_p > 2.4 \times 10^{34}$~yr ($p \to e^+ \pi^0$, 90\% CL)~\cite{SuperK2020}; the underlying CSc-1-allowed range from $G\mu$ would otherwise extend to $1.6 \times 10^{34}$~yr.
\end{table}

The most immediately testable prediction is the proton decay lifetime.
The CSc-1-allowed range of $G\mu$ maps directly to $E_\text{GUT}$ and hence to $\tau_p \propto E_\text{GUT}^4$, giving a raw window $[1.6, 5.5] \times 10^{34}$~yr before observational constraints.
The lower portion of this window, $\tau_p < 2.4 \times 10^{34}$~yr, is already excluded by Super-Kamiokande's $p \to e^+ \pi^0$ search~\cite{SuperK2020}, leaving the surviving prediction $\tau_p \in [2.4, 5.5] \times 10^{34}$~yr.
With hadronic matrix-element uncertainties of a factor $\sim 5$, this surviving window lies entirely within Hyper-Kamiokande's projected sensitivity.
\emph{If CSc-1 is confirmed as a cosmic string, Hyper-Kamiokande should observe proton decay early in its operating life}---or the SU(5) GUT interpretation is excluded.
The fact that Super-Kamiokande has not already seen the decay further sharpens the bound: the GUT scale lies at the upper end of the CSc-1-allowed band, which in turn predicts the higher end of the cosmic-string tension range, $G\mu/c^2 \approx 7 \times 10^{-7}$.

\subsection*{Falsification conditions}

The framework is designed to be falsifiable on multiple independent fronts.
The model is excluded if any of the following hold:
\begin{itemize}
\item Hyper-Kamiokande returns a null result for $p \to e^+ \pi^0$ across the full window $\tau_p \in [2.4, 5.5] \times 10^{34}$~yr.
\item The CSc-1 candidate is observationally disconfirmed---either by independent lensing follow-up showing the spectral pair is a chance projection, or by improved CMB and PTA bounds excluding $G\mu/c^2 \gtrsim 4 \times 10^{-7}$ for individual surviving strings.
\item The Totani $20$~GeV halo excess is shown to have an astrophysical origin (\emph{e.g.}, unresolved millisecond pulsar populations) inconsistent with a dark-matter annihilation morphology.
\item Direct-detection experiments report a positive signal at a mass scale incompatible with the dark proton ($\sim 187$~GeV) or any other element of the dark spectrum predicted by mode-locking with $\kappa_\text{GUT} = 5.32$.
\end{itemize}
Each of these channels is independent of the others, so a partial null result merely tightens the framework rather than refuting it.

The structural inaccessibility prediction is also sharp: dark matter \emph{cannot} be produced at any foreseeable collider, regardless of energy.
The barrier is not coupling strength but scale mismatch ($\xi_\text{GUT}/\xi_\text{SM} \sim 10^{-19}$).
Continued null results at the LHC and future colliders are a \emph{confirmation}, not a failure, of the two-condensate model.


%% ================================================================
\section{Discussion}
\label{sec:discussion}
%% ================================================================

We have proposed that string theory correctly identifies the fundamental objects of the dark sector: one-dimensional vortex lines in the $\psi_\text{GUT}$ condensate.
The visible sector, by contrast, is described by knot theory---compact vortex defects with aspect ratios of order unity.
Both descriptions emerge from the same superfluid vacuum physics, distinguished only by the aspect ratio $R/\xi$.

The framework is governed by two parameters: $\kappa_\text{SM} = \pi^2$ (a topological constant) and $\kappa_\text{GUT} = 5.32$ (determined by the GUT-scale physics).
From these two numbers---plus the universal constants $\hbar$, $c$, $G$, and $m_e$---the model aims to account for a broad set of particle-physics and cosmological observables: the Standard Model gauge group, 56 particle masses, the proposed dark particle spectrum, cosmic string tensions, the fine-structure constant, and the dark energy equation of state.

Within this proposal, string theory and NWT are presented as complementary rather than competing.
String theory supplies a mathematical language well suited to the dark sector, while NWT provides the underlying condensate picture in which the one-dimensional defects naturally live.
Taken together, the two frameworks aim to span an unusually wide range of scales, from the Planck length to the Hubble radius and from the electron mass to the cosmic string tension.

The concept dictionary (Table~\ref{tab:dictionary}) is an invitation.
Every entry represents a calculation that a string theorist could perform within the NWT framework, or an NWT prediction that string-theoretic methods could sharpen.
The Veneziano amplitude conjecture (\S\ref{sec:veneziano}) is the most specific such invitation: if the path integral over vortex reconnection configurations can be shown to yield the dual resonance model, the connection between NWT and string theory would be established at the level of the S-matrix.

\subsection{Relation to existing frameworks}

Several elements of the two-condensate model have precedent in the literature, and we wish to be explicit about what is inherited and what is new.

\textbf{Superfluid dark matter.}
The most important comparison is with the superfluid dark matter program of Berezhiani \& Khoury~\cite{BerezhianiKhoury2015}, in which dark matter forms a superfluid in galactic halos with phonon-mediated forces that reproduce MOND-like behavior.
The NWT condensate differs in a fundamental respect: its sound speed is $c_s = c$ (Lorentz-invariant vacuum), not $c_s \ll c$ (non-relativistic phonon).
Consequently, the NWT dark sector produces \emph{no} MOND-like phenomenology, no phonon-mediated forces, and no departure from CDM at any astrophysical scale above $\xi_\text{GUT}$.
The two models share the word ``superfluid'' but occupy opposite corners of parameter space.

\textbf{Fuzzy / ultralight dark matter.}
Ultralight axion models ($m \sim 10^{-22}$~eV) produce de~Broglie wavelength effects at kiloparsec scales~\cite{Hui2017}.
The NWT dark particles have masses of $\mathcal{O}(\text{GeV})$---forty orders of magnitude heavier---and their de~Broglie wavelengths are sub-femtometer.
The NWT dark sector is cold, collisionless dark matter at all observable scales, with no fuzzy-DM signatures.

\textbf{Hidden sectors and dark photons.}
The two-condensate model is, structurally, a hidden-sector model: a second gauge sector coupled to the Standard Model through a portal interaction.
This places it in a broad family that includes dark photon models, secluded sectors, and self-interacting hidden-sector dark matter.
The NWT-specific content is the identification of the hidden sector with the \emph{residual GUT condensate}, the derivation of the portal coupling from $\kappa_\text{SM} \times \kappa_\text{GUT}$ (Eq.~\ref{eq:alpha-dark}), and the prediction of a specific dark particle spectrum from mode-locking.
These are not features of generic hidden-sector models.

\textbf{Topological defects as dark matter.}
The idea that dark matter could consist of topological defects---cosmic string loops, domain wall fragments, or solitons---has a long history.
Our contribution is not the general concept but the specific realization: a quantitative mass spectrum from torus knot mode-locking, a derived interface coupling, and the identification of the visible/dark distinction with the aspect ratio of vortex defects ($R/\xi \sim 1$ versus $R/\xi \sim 10^{14}$).

\textbf{Cosmic strings from condensates.}
That superfluid condensates support string-like vortex defects is textbook physics.
The novelty here is the proposal that GUT-scale cosmic strings are vortex lines in a \emph{specific} residual condensate ($\psi_\text{GUT}$) whose network has largely decayed via Kelvin instability, leaving rare survivors (CSc-1) while converting most of its energy into dark matter particles (vortex rings).

\textbf{What is genuinely new.}
The novel contributions of this paper, to our knowledge without precedent, are:
(i)~the proposed string/knot duality---the identification of the visible and dark sectors with two aspect-ratio limits of the same vortex physics;
(ii)~the derivation of the interface coupling $\alpha_\text{dark} = 1/(\sqrt{2}\,\kappa_\text{SM}\kappa_\text{GUT})$;
(iii)~the dark particle spectrum from mode-locking with $\kappa_\text{GUT} = 5.32$;
(iv)~the unification of dark matter, dark energy, and cosmic strings as three stages of $\psi_\text{GUT}$ network evolution;
and (v)~the structural inaccessibility argument for why dark matter cannot be produced at colliders.

\subsection{Independent convergence}

Several recent results from independent groups converge on the NWT picture without invoking a superfluid vacuum.
Brandenberger et al.~\cite{Brandenberger2025} show that cosmic strings can resolve the JWST high-redshift galaxy abundance puzzle by seeding early structure formation---consistent with the $\psi_\text{GUT}$ network providing gravitational scaffolding at high $z$.
Allali, Jain \& Yan~\cite{Allali2025} simulate cosmic string loops captured by a Milky Way--like halo and find that loops with weak recoil closely trace the dark matter distribution---precisely the behavior expected if the loops \emph{are} the dark matter.
Slagter \& Pan~\cite{Slagter2017} show that cosmic string networks explain the observed alignment of quasar polarization axes on Mpc scales---another large-scale coherence signal consistent with a relic string network.

\section*{Acknowledgements}

We thank three AI reviewers---Gemini (Google), Perplexity, and ChatGPT (OpenAI)---for detailed referee-style critical reviews of earlier drafts of this manuscript.
Each shaped the final version in distinct ways.
Gemini identified the conflict between our originally stated proton-decay range and the existing Super-Kamiokande exclusion, leading to the narrower window $\tau_p \in [2.4, 5.5]\times 10^{34}$~yr now reported in \S\ref{sec:predictions} and the sharpening of the $G\mu/c^2 \approx 7\times 10^{-7}$ prediction for the cosmic-string tension.
Perplexity identified a tonal oscillation between calculation and conjecture in earlier drafts and supplied much of the line-by-line softening that informs the present epistemic posture, particularly in the abstract, the introduction, and the discussion of the candidate string-theoretic correspondences in \S\ref{sec:candidate-dualities}.
ChatGPT, playing the role of a strict round-two referee, flagged our use of the non-relativistic Gross--Pitaevskii equation as needing explicit framing as an analogue model rather than a fundamental theory, the need to delegate the crossing-algebra closure result to its proper home in~\cite{NWT8} and to label the SO(10) chirality-doubling extension as a conjectural Paper~10 proposal rather than an established result, and the over-strong language in the dictionary-table mapping of compactified dimensions to discrete quantum numbers.
All of these are addressed in the present version.
None of these contributors are responsible for any remaining errors.

%% ================================================================
%% BIBLIOGRAPHY
%% ================================================================

\begin{thebibliography}{99}

\bibitem{NWT1} J.~P.~Galasyn and C.~Th\'eodore, ``The Standard Model from a Torus Knot,'' Zenodo (2026), \href{https://doi.org/10.5281/zenodo.18891785}{10.5281/zenodo.18891785}.

\bibitem{NWT2} J.~P.~Galasyn and C.~Th\'eodore, ``Three Integers and a Mass,'' Zenodo (2026), \href{https://doi.org/10.5281/zenodo.18892311}{10.5281/zenodo.18892311}.

\bibitem{NWT3} J.~P.~Galasyn and C.~Th\'eodore, ``Nuclear Magic Numbers from Torus Topology,'' Zenodo (2026), \href{https://doi.org/10.5281/zenodo.19036783}{10.5281/zenodo.19036783}.

\bibitem{NWT4} J.~P.~Galasyn and C.~Th\'eodore, ``From Vortex Rings to the Top Quark,'' Zenodo (2026), \href{https://doi.org/10.5281/zenodo.19051589}{10.5281/zenodo.19051589}.

\bibitem{NWT5} J.~P.~Galasyn and C.~Th\'eodore, ``The Electron Mass from Phase Closure,'' Zenodo (2026), \href{https://doi.org/10.5281/zenodo.19072313}{10.5281/zenodo.19072313}.

\bibitem{NWT6} J.~P.~Galasyn and C.~Th\'eodore, ``The Particle Mass Spectrum from Torus Knot Mode-Locking,'' Zenodo (2026), \href{https://doi.org/10.5281/zenodo.19225259}{10.5281/zenodo.19225259}.

\bibitem{NWT7} J.~P.~Galasyn and C.~Th\'eodore, ``Charge, Leptons, and the Genus of Torus Knots,'' Zenodo (2026), \href{https://doi.org/10.5281/zenodo.19256133}{10.5281/zenodo.19256133}.

\bibitem{NWT8} J.~P.~Galasyn and C.~Th\'eodore, ``The Standard Model Gauge Group from Vortex Knot Topology,'' Zenodo (2026), \href{https://doi.org/10.5281/zenodo.19334925}{10.5281/zenodo.19334925}.

\bibitem{NWT9} J.~P.~Galasyn and C.~Th\'eodore, ``Coupling Constants, Mixing Matrices, and Cross Sections from Vortex Knot Surgery,'' Zenodo (2026), \href{https://doi.org/10.5281/zenodo.19335224}{10.5281/zenodo.19335224}.

\bibitem{Sazhina2019} O.~S.~Sazhina et al., ``Optical analysis of a CMB cosmic string candidate,'' Mon.\ Not.\ R.\ Astron.\ Soc.\ \textbf{485}, 1876 (2019).

\bibitem{Bulygin2023} I.~I.~Bulygin, M.~V.~Sazhin, and O.~S.~Sazhina, ``Theory of gravitational lensing on a curved cosmic string,'' Eur.\ Phys.\ J.\ C \textbf{83}, 844 (2023).

\bibitem{Sazhina2025} O.~S.~Sazhina, I.~I.~Bulygin, and O.~Y.~Solodilova, ``Problems and methods of modern search for cosmic strings,'' Astron.\ Rep.\ \textbf{69}, 14 (2025).

\bibitem{YusefZadeh2022} F.~Yusef-Zadeh et al., ``Statistical properties of the population of the Galactic Centre filaments,'' Astrophys.\ J.\ Lett.\ \textbf{925}, L18 (2022).

\bibitem{Totani2025} T.~Totani, ``20~GeV halo-like excess of the Galactic diffuse emission and implications for dark matter annihilation,'' J.\ Cosmol.\ Astropart.\ Phys.\ \textbf{2025}(11), 080 (2025); arXiv:2507.07209.

\bibitem{Cheng2026} H.~Cheng, E.~Di~Valentino, and L.~Visinelli, ``Cosmic strings as dynamical dark energy: novel constraints,'' J.\ High Energy Astrophys.\ (2026); doi:10.1016/j.jheap.2026.100610.

\bibitem{Swamy2025} I.~Swamy and D.~Singh, ``Blandford--Znajek jets and the total angular momentum evolution of a black hole connected to a cosmic string,'' Phys.\ Rev.\ D \textbf{112}, 043034 (2025).

\bibitem{Wang2021} P.~Wang, N.~I.~Libeskind, E.~Tempel et al., ``Possible observational evidence for cosmic filament spin,'' Nature Astron.\ \textbf{5}, 839 (2021).

\bibitem{Watkins2023} R.~Watkins, H.~A.~Feldman, and M.~J.~Hudson, ``Analysing the large-scale bulk flow using CosmicFlows4,'' Mon.\ Not.\ R.\ Astron.\ Soc.\ \textbf{524}, 1885 (2023).

\bibitem{JWST-Bullet2025} JWST Bullet Cluster lensing reconstruction, arXiv:2601.22245 (2025).

\bibitem{Kibble1976} T.~W.~B.~Kibble, ``Topology of cosmic domains and strings,'' J.\ Phys.\ A \textbf{9}, 1387 (1976).

\bibitem{Vilenkin1985} A.~Vilenkin, ``Cosmic strings and domain walls,'' Phys.\ Rep.\ \textbf{121}, 263 (1985).

\bibitem{Witten1985} E.~Witten, ``Superconducting strings,'' Nucl.\ Phys.\ B \textbf{249}, 557 (1985).

\bibitem{Polchinski2005} J.~Polchinski, ``Introduction to cosmic F- and D-strings,'' in \emph{String Theory: From Gauge Interactions to Cosmology}, NATO ASI (2005); hep-th/0412244.

\bibitem{Veneziano1968} G.~Veneziano, ``Construction of a crossing-symmetric, Regge-behaved amplitude for linearly rising trajectories,'' Nuovo Cim.\ A \textbf{57}, 190 (1968).

\bibitem{Maldacena1998} J.~M.~Maldacena, ``The large $N$ limit of superconformal field theories and supergravity,'' Adv.\ Theor.\ Math.\ Phys.\ \textbf{2}, 231 (1998).

\bibitem{Vollhardt1990} D.~Vollhardt and P.~W\"olfle, \emph{The Superfluid Phases of Helium 3} (Taylor \& Francis, 1990).

\bibitem{NANOGrav2023} G.~Agazie et al.\ (NANOGrav), ``The NANOGrav 15~yr data set: evidence for a gravitational-wave background,'' Astrophys.\ J.\ Lett.\ \textbf{951}, L8 (2023).

\bibitem{Ellis2021} J.~Ellis et al., ``Cosmic string interpretation of NANOGrav pulsar timing data,'' Phys.\ Rev.\ Lett.\ \textbf{126}, 041304 (2021).

\bibitem{Planck2013strings} P.~A.~R.~Ade et al.\ (Planck), ``Planck 2013 results.\ XXV.\ Searches for cosmic strings and other topological defects,'' Astron.\ Astrophys.\ \textbf{571}, A25 (2014).

\bibitem{BerezhianiKhoury2015} L.~Berezhiani and J.~Khoury, ``Theory of dark matter superfluidity,'' Phys.\ Rev.\ D \textbf{92}, 103510 (2015).

\bibitem{Hindmarsh2011} M.~B.~Hindmarsh, ``Signals of inflationary models with cosmic strings,'' Prog.\ Theor.\ Phys.\ Suppl.\ \textbf{190}, 197 (2011).

\bibitem{Hui2017} L.~Hui, J.~P.~Ostriker, S.~Tremaine, and E.~Witten, ``Ultralight scalars as cosmological dark matter,'' Phys.\ Rev.\ D \textbf{95}, 043541 (2017).

\bibitem{ACT-DR6-strings} A.~Avgoustidis, E.~J.~Copeland, A.~Moss, and J.~Raidal, ``CMB anisotropies from cosmic (super)strings in light of ACT DR6,'' arXiv:2602.18272 (2026).

\bibitem{YusefZadeh2025JWST} F.~Yusef-Zadeh et al., ``Non-stop variability of Sgr~A* using JWST at 2.1 and 4.8~$\mu$m wavelengths,'' arXiv:2501.04096 (2025).

\bibitem{YusefZadeh2022spacing} F.~Yusef-Zadeh, R.~G.~Arendt, M.~Wardle et al., ``Statistical properties of the population of the Galactic Centre filaments~II: the spacing between filaments,'' Mon.\ Not.\ R.\ Astron.\ Soc.\ \textbf{515}, 3059 (2022).

\bibitem{YusefZadeh2023horizontal} F.~Yusef-Zadeh et al., ``The population of the Galactic Center filaments: position angle distribution reveals a degree-scale collimated outflow from Sgr~A* along the Galactic plane,'' Astrophys.\ J.\ Lett.\ \textbf{949}, L31 (2023).

\bibitem{Swamy2025accretion} I.~Swamy and D.~Singh, ``Investigating the interaction of a cosmic string with an accreting black hole,'' arXiv:2512.09380 (2025).

\bibitem{Swamy2025QNM} I.~Swamy and D.~Singh, ``Probing cosmic strings via black hole quasinormal modes in gravitational wave astronomy,'' arXiv:2512.08368 (2025).

\bibitem{Swamy2025Xray} I.~Swamy and D.~Singh, ``Impact on orbital period of X-ray binary system attached to a cosmic string,'' arXiv:2512.07911 (2025).

\bibitem{Swamy2025DM} I.~Swamy and D.~Singh, ``Genesis of axion dark matter via black hole and cosmic string dynamics,'' arXiv:2512.07268 (2025).

\bibitem{Tranchedone2026} L.~Tranchedone, E.~Carragher, E.~Hardy, and N.~Koscelansk\'a van IJcken, ``Metastable cosmic strings are broken at the start,'' arXiv:2601.04320 (2026).

\bibitem{Penna2025} R.~Penna, ``Semiclassical cosmic string evaporation,'' arXiv:2510.18948 (2025).

\bibitem{Auclair2025} P.~Auclair, ``Breaking the strings: the signatures of cosmic string loop fragmentation,'' J.\ Cosmol.\ Astropart.\ Phys.\ \textbf{2025}(12), 020 (2025).

\bibitem{Brandenberger2025} M.~Blamart, A.~Liu, R.~Brandenberger, J.~B.~Mu\~noz, and B.~Cyr, ``UV luminosity functions from HST and JWST: a possible resolution to the high-redshift galaxy abundance puzzle and implications for cosmic strings,'' arXiv:2512.09980 (2025).

\bibitem{Allali2025} I.~Allali, M.~Jain, and S.~Yan, ``Clustering of cosmic string loops within a Milky-Way-like halo,'' arXiv:2512.23811 (2025).

\bibitem{Slagter2017} R.~J.~Slagter and S.~Pan, ``Evidence of cosmic strings by the observation of the alignment of quasar polarization axes,'' Zenodo (2017), doi:10.5281/zenodo.1038103.

\bibitem{Meichanetzidis2025} K.~Meichanetzidis et al., ``Quantum advantage and CSS-T codes from Jones polynomials,'' arXiv:2503.05625 (2025).

\bibitem{NWTJonesQC} J.~P.~Galasyn and C.~Th\'eodore, ``Jones polynomials of the NWT particle knots on superconducting quantum hardware,'' (in preparation, 2026).

\bibitem{SuperK2020} A.~Takenaka et al.\ (Super-Kamiokande), ``Search for proton decay via $p \to e^+\pi^0$ and $p \to \mu^+\pi^0$ with an enlarged fiducial volume of Super-Kamiokande~I--IV,'' Phys.\ Rev.\ D \textbf{102}, 112011 (2020).

\end{thebibliography}

\section*{Data and Code Availability}

All analysis code, simulation scripts, and the LaTeX source for this manuscript are available at \url{https://github.com/JimGalasyn/null-worldtube}.

\end{document}
